\chapter{\Glsentryshort{milp} Representation of \Glsentryshort{bra} Routing Modes}
\label{ap:math}

BRA exposes two routing modes: a scenario-based candidate-selection model and an \gls{lbh} mode. The scenario mode is naturally represented as a \gls{milp} over a generated set of candidate routes. The \gls{lbh} mode is implemented as a constructive insertion heuristic, but it can be interpreted as a greedy approximation to a related bus-itinerary assignment problem.

\section{Scenario Candidate-Selection Formulation}

Let $M$ denote the set of students, $S$ the set of schools, $B$ the set of buses, and $R$ the set of candidate single-school routes generated by the scenario procedure. Each route $r\in R$ serves a subset of students $M_r \subseteq M$, visits a school $s\left(r\right)\in S$, has travel cost $c_r$, travel time $t_r$, and requires a vehicle type compatible with the students assigned to it. Let $a_{m,r}=1$ if student $m$ is served by candidate route $r$, and 0 otherwise. Let $K\subseteq R\times R$ be the set of feasible route links, where $\left(r,r'\right)\in K$ means that the same bus can complete route $r$ and then route $r'$ while respecting school arrival times, deadhead travel time, and required buffer times.

We introduce binary variables
\[
x_r =
\begin{cases}
1 & \text{if route } r \text{ is selected},\\
0 & \text{otherwise,}
\end{cases}
\]
\[
z_b =
\begin{cases}
1 & \text{if bus } b \text{ is used},\\
0 & \text{otherwise,}
\end{cases}
\]
\[
y_{b,r} =
\begin{cases}
1 & \text{if route } r \text{ is assigned to bus } b,\\
0 & \text{otherwise,}
\end{cases}
\]
and
\[
\ell_{b,rr'} =
\begin{cases}
1 & \text{if bus } b \text{ operates route } r' \text{ immediately after route } r,\\
0 & \text{otherwise.}
\end{cases}
\]

A simplified scenario formulation is:
\[
\min
\sum_{r\in R} c_r x_r
+ \lambda_B \sum_{b\in B} z_b
+ \lambda_T \sum_{r\in R} t_r x_r
+ \lambda_U \sum_{m\in M} u_m ,
\]
subject to
\[
\sum_{r\in R} a_{m,r}x_r + u_m = 1
\qquad \forall m\in M,
\]
\[
\sum_{b\in B} y_{b,r} = x_r
\qquad \forall r\in R,
\]
\[
y_{b,r} \le z_b
\qquad \forall b\in B,\ r\in R,
\]
\[
y_{b,r} = 0
\qquad \forall b\in B,\ r\in R \text{ such that } b \text{ is incompatible with } r,
\]
\[
\sum_{r\in R} y_{b,r} \le Q_{\max} z_b
\qquad \forall b\in B,
\]
\[
\ell_{b,rr'} \le y_{b,r},
\qquad
\ell_{b,rr'} \le y_{b,r'}
\qquad \forall b\in B,\ \left(r,r'\right)\in K,
\]
\[
\ell_{b,rr'} = 0
\qquad \forall b\in B,\ \left(r,r'\right)\notin K,
\]
\[
\sum_{r':\left(r,r'\right)\in K} \ell_{b,rr'} \le y_{b,r}
\qquad \forall b\in B,\ r\in R,
\]
\[
\sum_{r':\left(r',r\right)\in K} \ell_{b,r'r} \le y_{b,r}
\qquad \forall b\in B,\ r\in R,
\]
\[
x_r,y_{b,r},z_b,\ell_{b,rr'}\in\left\{0,1\right\},
\qquad
u_m\in\left\{0,1\right\}.
\]

Here, $u_m$ is an optional unserved-student variable. Setting $\lambda_U$ very large enforces near-full coverage whenever feasible. The generated route set $R$ already encodes many local feasibility requirements, including capacity, school compatibility, maximum ride time, and wheelchair or monitor requirements. The \gls{milp} therefore optimizes over feasible route candidates rather than over every street-level arc.

\section{\Glsentryshort{lbh} as a Greedy Approximation to an Itinerary \Glsentryshort{milp}}

The \gls{lbh} mode is not solved as a \gls{milp}. It builds bus itineraries directly, processing buses in capacity priority order and inserting uncovered stops into the currently constructed itinerary whenever the insertion is feasible. However, the heuristic can be interpreted as approximating the following itinerary-selection problem.

Let $I_b$ be the set of feasible itineraries available to bus $b$. An itinerary $i\in I_b$ is an ordered sequence of one or more school routes that can be operated by bus $b$. Let $a_{m,b,i}=1$ if itinerary $i$ for bus $b$ serves student $m$, and let $c_{b,i}$ be the cost of operating that itinerary. Define
\[
w_{b,i} =
\begin{cases}
1 & \text{if bus } b \text{ operates itinerary } i,\\
0 & \text{otherwise.}
\end{cases}
\]

The corresponding set-packing formulation is:
\[
\min
\sum_{b\in B}\sum_{i\in I_b} c_{b,i}w_{b,i}
+ \lambda_B \sum_{b\in B} z_b
+ \lambda_U \sum_{m\in M} u_m ,
\]
subject to
\[
\sum_{b\in B}\sum_{i\in I_b} a_{m,b,i}w_{b,i} + u_m = 1
\qquad \forall m\in M,
\]
\[
\sum_{i\in I_b} w_{b,i} \le z_b
\qquad \forall b\in B,
\]
\[
w_{b,i}\in\left\{0,1\right\},
\qquad
z_b\in\left\{0,1\right\},
\qquad
u_m\in\left\{0,1\right\}.
\]

The difficulty is that \(I_b\) is enormous: each itinerary corresponds to many possible ordered insertions of stops and schools. \gls{lbh} avoids explicitly enumerating or optimizing over this full set. Instead, it constructs one itinerary at a time using cheapest feasible insertion. At each step, for the current partial itinerary \(i\), the heuristic chooses an uncovered stop or student group \(g\) that minimizes the marginal insertion cost
\[
\Delta c\left(i,g\right)
\]
subject to feasibility:
\[
\operatorname{cap}\left(i\oplus g\right) \le C_b,
\]
\[
\operatorname{ride}_m\left(i\oplus g\right) \le H_m
\qquad \forall m \text{ served by } i\oplus g,
\]
\[
\operatorname{arr}_{s}\left(i\oplus g\right) \in \left[\underline{h}_s,\overline{h}_s\right]
\qquad \forall s \text{ visited by } i\oplus g,
\]
\[
\operatorname{type}\left(b\right) \text{ is compatible with all students served by } i\oplus g.
\]

Thus, the \gls{lbh} decision rule can be written as
\[
g^\star
=
\arg\min_{g\in U}
\left\{
\Delta c\left(i,g\right)
:
i\oplus g \begin{tabular}{l} satisfies capacity, timing, ride-time, \\ grade, and accessibility constraints\end{tabular} %\text{ satisfies capacity, timing, ride-time, grade, and accessibility constraints}
\right\},
\]
where $U$ is the set of currently uncovered students or stops. The algorithm repeats this insertion step until no feasible insertion remains, then proceeds to the next bus.

In this sense, the scenario mode solves a global route-selection MILP over generated candidates, while \gls{lbh} greedily constructs a feasible solution to an implicit itinerary-selection problem without solving the full \gls{milp}.
